\documentclass[11pt,a4paper]{article}

% ---------- Packages ----------
\usepackage[utf8]{inputenc}
\usepackage[T1]{fontenc}
\usepackage{lmodern}
\usepackage[english]{babel}
\usepackage{geometry}
\geometry{a4paper, margin=1in}
\usepackage{graphicx}
\usepackage{hyperref}
\usepackage{xcolor}
\usepackage{listings}
\usepackage{booktabs}
\usepackage{tabularx}
\usepackage{caption}
\usepackage{enumitem}
\usepackage{parskip}
\usepackage{float}
\usepackage[backend=biber,style=numeric-comp,sorting=none]{biblatex}
\addbibresource{refs.bib}

% ---------- Hyperref styling ----------
\hypersetup{
    colorlinks=true,
    linkcolor=black,
    urlcolor=blue,
    citecolor=blue
}

% ---------- Listings (code) ----------
\definecolor{codegray}{gray}{0.95}
\lstset{
    backgroundcolor=\color{codegray},
    basicstyle=\ttfamily\small,
    breaklines=true,
    frame=single,
    captionpos=b,
    numbers=left,
    numberstyle=\tiny\color{gray},
    showstringspaces=false,
    tabsize=2
}

% ---------- Title ----------
\title{
    \textbf{Data Warehousing Strategies for an E-commerce Domain} \\
    \large IKT553 -- Intelligent Database Management
}
\author{Group Project Report}
\date{\today}

\begin{document}

\maketitle

% =====================================================
\section*{Acknowledgements}
\addcontentsline{toc}{section}{Acknowledgements}

We would like to thank our course instructor and teaching assistants for their
guidance throughout this project. We also thank the open-source community behind
PostgreSQL, MongoDB, Neo4j, Apache Kafka, ksqlDB, Flask and Docker, whose tools
made the hands-on parts of this work possible.

% =====================================================
\section*{Abstract}
\addcontentsline{toc}{section}{Abstract}

This report presents the design and implementation of a Data Warehouse (DW)
for an e-commerce domain. The goal was to support common business questions
such as ``which products generate the most revenue?'' or ``how does profit
change over time?''. We built the same logical warehouse in three different
technologies: a relational star schema in PostgreSQL, a document model in
MongoDB, and a graph model in Neo4j. On top of these, we built a single
Flask-based ``umbrella'' web application that lets a user run the same
decision-support query on any of the three backends and compare the results
side by side. The whole stack runs in Docker, and query executions are
streamed through Kafka and aggregated with ksqlDB to give a near real-time
view of how the warehouse is being used. The report describes the data model,
the loading process, the front-end, the streaming pipeline, and a comparison
of the three database technologies from a data-warehousing point of view.

% =====================================================
\tableofcontents
\listoffigures
\listoftables
\newpage

% =====================================================
\section{Introduction}

E-commerce platforms produce a lot of transactional data: orders, customers,
products, shipping events, and returns. On its own, this raw data is useful
for running the store, but it is not shaped in a way that makes it easy to
answer business questions. A Data Warehouse reorganises this information
around the decisions that managers and analysts actually need to make, so
queries like ``which region gives us the highest profit?'' can be answered
quickly and consistently.

In this project we build a data warehouse for a simplified e-commerce
business and then implement it three times: once in a relational database
(PostgreSQL), once in a document database (MongoDB), and once in a graph
database (Neo4j). A single web front-end lets us run queries against any of
them and compare the results. We also add a streaming layer with Kafka and
ksqlDB to observe how queries are being used over time.

\subsection{Report Structure}

The report is organised as follows. Section~\ref{sec:background} introduces
the main technologies we used. Section~\ref{sec:methods} explains the design
choices: the star schema, the MongoDB document shape, the Neo4j graph model,
and the architecture of the web application and streaming layer.
Section~\ref{sec:results} shows what the system looks like in action: the
front-end, the same query running on all three backends, and the
ksqlDB-derived metrics. Section~\ref{sec:discussion} compares the three
implementations and discusses trade-offs. Section~\ref{sec:conclusion}
concludes and Section~\ref{sec:ai} describes how AI tools were used during
the project.

\subsection{Scope and Motive}

The scope of the project is a single e-commerce domain with sales, products,
customers, locations, shipping modes, and dates as the main business
entities. We focus on decision-support queries (aggregations, rankings,
trends) rather than transactional workloads. The motive is practical: we
wanted to understand, by building it, how classic data-warehousing concepts
such as the star schema translate into NoSQL systems that are structured
very differently.

\subsection{Assumptions}

To keep the scope manageable, we made the following assumptions:

\begin{itemize}[leftmargin=*]
    \item The data is treated as already-cleaned business data. The project
          focuses on modelling and querying, not on building a full ETL
          pipeline from messy source systems.
    \item The three backends contain the same semantic data. We do not
          attempt to keep them in real-time sync; they are all loaded from
          the same source material.
    \item The front-end targets an analyst or decision-maker, not the
          general public, so authentication and per-user access control
          are out of scope.
    \item Kafka and ksqlDB are used for \emph{observability} of query
          executions (which queries ran, how fast, on which DB), not for
          streaming business events into the warehouse.
    \item All components run locally through Docker. We assume a developer
          machine with Docker Desktop and at least a few gigabytes of free
          memory.
\end{itemize}

% =====================================================
\section{Background}\label{sec:background}

This section introduces the core ideas and technologies that the rest of
the report builds on. We cover the two database families (relational and
NoSQL), the data-warehousing concepts we apply to each, and the supporting
tools that let us package, run, and observe the whole system.

\subsection{Relational Databases and PostgreSQL}

The relational model, introduced by Codd in 1970~\cite{codd1970relational},
organises data into tables of tuples linked by shared keys. It became
dominant because it separates the logical shape of the data from how it is
stored and gives users a declarative query language. Over the following
decades the model was extended with transactions, indexing, cost-based
query optimisation, and specialised analytical operators, while remaining
the backbone of most enterprise systems~\cite{stonebraker2010sqlvsnosql}.

PostgreSQL is a mature, open-source object-relational database management
system that supports standard SQL, ACID transactions, JSON data types,
window functions and materialised views, all of which are useful for
analytical workloads~\cite{postgresql_docs}. In our project PostgreSQL
plays the role of the \emph{relational data warehouse}: the star schema
described in Section~\ref{sec:star-schema} is created there, and we use
Supabase~\cite{supabase_docs} as the hosted PostgreSQL environment so the
warehouse is reachable without running an extra database container.

\subsection{NoSQL Databases and the Document Model (MongoDB)}

``NoSQL'' is an umbrella term for database systems that step away from
strict relational tables to scale horizontally, tolerate flexible schemas,
or match a specific data shape more naturally. Early motivations came from
large web services such as Amazon's Dynamo~\cite{decandia2007dynamo},
which prioritised availability and horizontal scalability over strong
consistency. Surveys by Cattell~\cite{cattell2011nosql},
Leavitt~\cite{leavitt2010nosql}, Moniruzzaman and
Hossain~\cite{moniruzzaman2013nosql}, and later Gessert et
al.~\cite{gessert2017nosql} divide modern NoSQL systems into four main
families -- key-value, document, wide-column, and graph stores -- and
discuss their trade-offs in detail. The CAP-based intuition that
distributed databases must choose between availability and consistency
during a partition is further refined by Abadi's PACELC
formulation~\cite{abadi2012pacelc}, which also considers the
latency/consistency trade-off in normal operation. Our project uses one
document store (MongoDB) and one graph database (Neo4j).

MongoDB is a widely used document database that stores data as BSON
documents -- essentially JSON with a richer type system~\cite{mongodb_docs}.
Instead of splitting related information across several normalised tables,
MongoDB encourages \emph{embedding}: a single document can contain its own
nested sub-documents and arrays. In a data-warehouse context this matters
because a sale fact can be stored together with its product, customer,
location and date information in one document, which removes the need for
joins at query time. Queries are written as \emph{aggregation
pipelines}~\cite{mongodb_aggregation}: a sequence of stages
(\texttt{\$match}, \texttt{\$group}, \texttt{\$sort}, \texttt{\$project})
that transform documents into the final result.

\subsection{Graph Databases and Neo4j}

Graph databases treat relationships as first-class citizens: data is
modelled as nodes, labelled edges, and properties on
both~\cite{angles2008graph,robinson2015graph}. This makes them especially
strong for questions that involve multi-hop connections, paths, or
patterns between entities, which are expensive to express as repeated
joins in a relational system.

Neo4j is one of the most mature graph databases and uses the Cypher query
language, an ASCII-art-like syntax where the query visually mirrors the
pattern being matched~\cite{neo4j_docs,cypher_docs,francis2018cypher}. For
example, \texttt{(s:Sale)-[:OF\_PRODUCT]->(p:Product)} reads literally as
``a sale that is of a product''. In this project Neo4j plays the role of
our second NoSQL implementation: it stores the same business facts as the
relational warehouse, but reorganised so that dimensions (product,
customer, location, date, ship mode) become labelled nodes and each sale
becomes a node connected to them through typed relationships.

\subsection{Data Warehouse}

A \emph{data warehouse} is a database optimised for analytical decision
support rather than for day-to-day transactions. Chaudhuri and Dayal give
the classic overview: data is extracted from one or more operational
source systems, cleaned, integrated around business subjects (customers,
sales, products), and stored in a form that is easy to aggregate over
long time ranges~\cite{chaudhuri1997olap}.

Data-warehouse workloads differ from transactional (OLTP) workloads in
three important ways~\cite{chaudhuri1997olap,kimball2013toolkit}: queries
tend to be read-heavy and touch many rows, the schema is tuned for fast
aggregation rather than fast insertion, and the data is loaded in batches
(an ETL or ELT process) rather than row by row. Cloud-scale variants of
this idea are surveyed by Grolinger et al.~\cite{grolinger2013cloud}.
Understanding this distinction is important for our project because it
drives all of the modelling choices that follow.

\subsection{Star Schema}

The star schema~\cite{kimball2013toolkit} is the most common way to
structure a relational data warehouse. It has one central \emph{fact
table} that stores the measurements of the business process (for us: sales
amount, profit amount, quantity) and several surrounding \emph{dimension
tables} that describe the context (who bought it, what was sold, where it
went, when, and how). Each dimension is linked to the fact table through a
foreign key, which gives the schema its characteristic star shape.

Its popularity comes from three properties. First, it is \emph{easy for
humans to reason about}: most business questions translate directly into
``slice the facts by these dimensions''. Second, it is \emph{efficient to
query}: a single join from fact to dimension is cheap, and modern
optimisers handle star-join queries well. Third, it is \emph{BI-tool
friendly}: reporting tools assume this layout and can generate queries
automatically~\cite{kimball2013toolkit}. An important part of our project
is to see how far this idea travels when we leave the relational world:
what a ``fact'' and a ``dimension'' look like in a document database and in
a graph database.

\subsection{Docker}

Docker is a container platform that packages an application together with
its dependencies into an immutable image that can be run reproducibly on
any Linux-compatible host~\cite{docker_docs,merkel2014docker}. Unlike a
full virtual machine, a container shares the host's kernel and is
therefore much lighter to start and to distribute.

For a project like ours, which uses three different databases, a streaming
platform, a schema registry and ksqlDB, the practical value of Docker is
immense: a single \texttt{docker-compose.yml} file describes all of these
services, their network links, their ports, and their environment, and one
command (\texttt{docker compose up}) brings the whole stack
up~\cite{docker_compose_docs}. This makes the project reproducible across
machines and removes the ``works on my laptop'' problem.

\subsection{Apache Kafka and ksqlDB}

Apache Kafka is a distributed log-based streaming platform originally
developed at LinkedIn~\cite{kreps2011kafka,kafka_docs}. Producers publish
messages to named \emph{topics}, consumers subscribe to those topics, and
Kafka stores the messages in an append-only log that can be replayed. The
``log as a unifying abstraction'' idea is articulated in Kreps' widely
cited essay~\cite{kreps2013log}. Because the log decouples producers and
consumers in time, Kafka is widely used to build asynchronous pipelines
between services.

ksqlDB~\cite{ksqldb_docs} is an event-streaming database built on top of
Kafka. It is grounded in the stream/table duality formalised by Sax et
al.~\cite{sax2018streams,confluent_streams_tables} and lets developers
express stream processing in a SQL-like language: define a \texttt{STREAM}
over a Kafka topic, then derive new streams and tables from it using
familiar constructs such as \texttt{SELECT}, \texttt{WHERE}, \texttt{GROUP
BY} and windowed aggregations. In this project we do not put business
data through Kafka; instead we use Kafka + ksqlDB to observe the
warehouse \emph{itself}: every query executed through the web UI produces
a small event, and ksqlDB turns that event stream into rolling metrics
about query latency, volume and failures.

\subsection{Python}

Python is a high-level, dynamically typed programming language that is
widely used for data-oriented work and for building small web
services~\cite{python_docs}. We chose it as the implementation language
for the back-end for two reasons. First, it has well-maintained client
libraries for every component in our stack: \texttt{psycopg2} for
PostgreSQL~\cite{psycopg_docs}, \texttt{pymongo} for
MongoDB~\cite{pymongo_docs}, the official Python driver for
Neo4j~\cite{neo4j_python_docs}, and \texttt{kafka-python} for the Kafka
producer~\cite{kafka_python_docs}. Second, its syntax makes it easy to
read by anyone on the team regardless of background.

\subsection{Flask}

Flask is a lightweight Python web framework based on the Werkzeug WSGI
toolkit and the Jinja2 template engine~\cite{flask_docs}. It is often
described as a \emph{microframework} because its core keeps very few
assumptions about how an application should be structured; extra features
are added through extensions as needed.

For a single-developer or small-team project like ours, these properties
fit well: we only need a handful of HTML pages, two JSON endpoints
(\texttt{/api/run} and \texttt{/api/compare}), and a thin glue layer that
dispatches a query to the correct database client and publishes an event
to Kafka. Flask gives us exactly that without any ceremony.

% =====================================================
\section{Methods}\label{sec:methods}

This section describes what we actually built. It covers the data model in
each database, how the three implementations relate to each other, and how
the front-end and the streaming layer tie everything together.

\subsection{System Overview}

The system has three kinds of users in mind:

\begin{itemize}[leftmargin=*]
    \item \textbf{Business analysts}, who want to answer decision-support
          questions using predefined queries and the side-by-side comparison
          view.
    \item \textbf{Technical users/administrators}, who want to write custom
          queries, inspect query performance, or check the health of the
          streaming layer through ksqlDB.
    \item \textbf{Developers}, who extend the set of predefined queries or
          change the data model.
\end{itemize}

All users interact with the same Flask application. There is no separate UI
per database: the ``umbrella'' front-end hides the underlying technology and
offers the same experience regardless of which backend is selected.

\begin{figure}[H]
\centering
% TODO: High-level architecture diagram. Boxes for:
%  - Browser (the umbrella UI)
%  - Flask app (umbrella)
%  - PostgreSQL (Supabase), MongoDB, Neo4j
%  - Kafka broker + topic "dw.query.executions"
%  - ksqlDB server with QUERY_EXECUTIONS_RAW / QUERY_EXECUTION_METRICS /
%    FAILED_QUERY_EXECUTIONS
% Arrows: browser -> Flask (HTTP); Flask -> 3 DBs (native drivers);
%         Flask -> Kafka (producer); Kafka -> ksqlDB (stream processing).
\fbox{\parbox{0.9\textwidth}{\centering
\textit{Placeholder: system architecture diagram showing the browser,
the Flask umbrella app, the three databases, Kafka, and ksqlDB with
its derived streams/tables.}
\vspace{5em}}}
\caption{High-level architecture of the umbrella data warehouse.}
\label{fig:architecture}
\end{figure}

\subsection{Data Description}

The business domain is a generic e-commerce store. The facts are
\emph{sales}: each row represents a line item from an order. Each sale has
measurable values (sales amount, profit amount, quantity) and links to
several descriptive entities:

\begin{itemize}[leftmargin=*]
    \item The \textbf{product} that was sold, together with its category.
    \item The \textbf{customer} who bought it.
    \item The \textbf{location} the order was shipped to (region, state).
    \item The \textbf{shipping mode} used.
    \item The \textbf{date} of the transaction, broken down into year,
          month, and month name for trend queries.
\end{itemize}

These are exactly the entities that become tables in the relational model,
embedded sub-documents in MongoDB, and nodes in Neo4j.

\subsection{Data Loading Process}

The loading process follows a simple ``source $\rightarrow$ clean $\rightarrow$
shape'' pattern.

\begin{enumerate}[leftmargin=*]
    \item Raw e-commerce data (CSV-like records of orders and products) is
          taken as the source.
    \item Light cleaning is applied: trimming whitespace, normalising date
          formats, resolving product categories, and dropping rows with
          missing critical fields.
    \item The cleaned data is then shaped differently for each database:
          flattened into fact/dimension rows for PostgreSQL, bundled into
          sales documents for MongoDB, and split into nodes and
          relationships for Neo4j.
\end{enumerate}

Because the assignment scope focuses on modelling rather than on a full
production ETL, we treat the cleaned dataset as our input and describe the
transformations at a specification level.

\subsection{Star Schema Design}\label{sec:star-schema}

Following Kimball's dimensional modelling
guidelines~\cite{kimball2013toolkit}, the relational warehouse is
organised as a star schema with one fact table and five dimension tables.
Figure~\ref{fig:star-schema} shows the actual layout as it exists in our
Supabase-hosted PostgreSQL instance.

\begin{figure}[H]
\centering
% TODO: Insert the Supabase ERD screenshot for the star schema.
% Should show fact_sales at the centre linked by dashed lines to
% dim_date, dim_product, dim_customer, dim_location, dim_ship_mode.
\fbox{\parbox{0.9\textwidth}{\centering
\textit{Placeholder: ERD screenshot from Supabase showing
\texttt{fact\_sales} at the centre, linked to \texttt{dim\_date},
\texttt{dim\_product}, \texttt{dim\_customer}, \texttt{dim\_location},
and \texttt{dim\_ship\_mode}.}
\vspace{7em}}}
\caption{Star schema for the e-commerce data warehouse (PostgreSQL / Supabase).}
\label{fig:star-schema}
\end{figure}

\textbf{Fact table:} \texttt{fact\_sales}, with the measures
\texttt{sales\_amount}, \texttt{profit\_amount} and \texttt{quantity}, plus
foreign keys to each dimension (\texttt{product\_key},
\texttt{customer\_key}, \texttt{location\_key}, \texttt{date\_key},
\texttt{ship\_mode\_key}) and a surrogate primary key
\texttt{sales\_fact\_id}.

\textbf{Dimension tables:}
\begin{itemize}[leftmargin=*]
    \item \texttt{dim\_product} -- product name, category name.
    \item \texttt{dim\_customer} -- customer name (and could be extended
          with segment, address, etc.).
    \item \texttt{dim\_location} -- region, state.
    \item \texttt{dim\_date} -- year, month, month name.
    \item \texttt{dim\_ship\_mode} -- shipping mode name.
\end{itemize}

A simplified DDL outline is shown in Listing~\ref{lst:ddl}. (In the real
project, each dimension also has a surrogate key column.)

\begin{lstlisting}[language=SQL, caption={DDL outline for the star schema}, label={lst:ddl}]
CREATE TABLE dim_product (
  product_key    SERIAL PRIMARY KEY,
  product_name   TEXT NOT NULL,
  category_name  TEXT
);

CREATE TABLE dim_customer (
  customer_key   SERIAL PRIMARY KEY,
  customer_name  TEXT NOT NULL
);

CREATE TABLE dim_location (
  location_key   SERIAL PRIMARY KEY,
  region         TEXT,
  state          TEXT
);

CREATE TABLE dim_date (
  date_key       SERIAL PRIMARY KEY,
  year           INT,
  month          INT,
  month_name     TEXT
);

CREATE TABLE dim_ship_mode (
  ship_mode_key    SERIAL PRIMARY KEY,
  ship_mode_name   TEXT
);

CREATE TABLE fact_sales (
  sales_fact_id   SERIAL PRIMARY KEY,
  date_key        INT REFERENCES dim_date(date_key),
  product_key     INT REFERENCES dim_product(product_key),
  customer_key    INT REFERENCES dim_customer(customer_key),
  location_key    INT REFERENCES dim_location(location_key),
  ship_mode_key   INT REFERENCES dim_ship_mode(ship_mode_key),
  sales_amount    NUMERIC(12,2),
  quantity        INT,
  profit_amount   NUMERIC(12,2)
);
\end{lstlisting}

\subsubsection*{Pre-Aggregated Summary Tables}

Even with indexes, running the same aggregation queries over every single
sale becomes slow as the warehouse grows. To address this, we define
pre-aggregated summary tables that are refreshed by a batch job. Two
examples are shown in Listing~\ref{lst:summary}.

\begin{lstlisting}[language=SQL, caption={Example summary tables}, label={lst:summary}]
CREATE TABLE summary_monthly_revenue AS
SELECT d.year, d.month, SUM(f.sales_amount) AS revenue
FROM fact_sales f JOIN dim_date d ON f.date_key = d.date_key
GROUP BY d.year, d.month;

CREATE TABLE summary_category_revenue AS
SELECT p.category_name, SUM(f.sales_amount) AS total_sales
FROM fact_sales f JOIN dim_product p ON f.product_key = p.product_key
GROUP BY p.category_name;
\end{lstlisting}

These summary tables are separate from the nightly batch job that loads
\texttt{fact\_sales} and the dimensions. They are refreshed afterwards, in
their own batch step, and only the SQL specification is provided here.

\subsection{MongoDB (Document) Implementation}

MongoDB does not have the concept of ``tables joined by keys''. Instead, we
store a \emph{single document per sale} in a \texttt{sales} collection, with
the dimension information embedded as sub-documents:

\begin{lstlisting}[language=Python, caption={Shape of a MongoDB sales document (simplified)}]
{
  "_id": ObjectId("..."),
  "sales_amount": 199.99,
  "quantity": 2,
  "profit_amount": 34.00,
  "product": { "product_name": "Headphones", "category": "Electronics" },
  "customer": { "customer_name": "Alice Smith" },
  "location": { "region": "West", "state": "California" },
  "ship_mode": { "ship_mode_name": "Standard Class" },
  "date": { "year": 2024, "month": 3, "month_name": "March" }
}
\end{lstlisting}

The advantage is that all the information needed to answer a sales question
is already in the document, so there are no joins. The disadvantage is that
dimension information is duplicated across documents and becomes harder to
update consistently.

Queries are expressed as aggregation pipelines. For example, the ``top 10
products by revenue'' query groups by \texttt{product.product\_name} and
sums \texttt{sales\_amount}.

\subsection{Neo4j (Graph) Implementation}

In Neo4j, everything becomes a node or a relationship. The nodes we use
are: \texttt{Sale}, \texttt{Product}, \texttt{Category}, \texttt{Customer},
\texttt{State}, \texttt{ShipMode}, \texttt{Date}, and \texttt{Month}. The
relationships connect a \texttt{Sale} to its product
(\texttt{OF\_PRODUCT}), its customer (\texttt{BY\_CUSTOMER}), its
destination (\texttt{SHIPPED\_TO}), its shipping mode
(\texttt{USES\_SHIP\_MODE}), and its date (\texttt{ON\_DATE}).

Queries are written in Cypher and follow a ``pattern matching'' style:

\begin{lstlisting}[caption={Cypher: top 10 products by revenue}]
MATCH (s:Sale)-[:OF_PRODUCT]->(p:Product)
RETURN p.product_name AS product_name,
       SUM(s.sales_amount) AS revenue
ORDER BY revenue DESC
LIMIT 10;
\end{lstlisting}

The shape of the query reflects the shape of the data: ``starting from a
sale, follow the OF\_PRODUCT relationship to a product, and aggregate''.

\subsection{Queries and Front-End}

The system ships with ten pre-defined decision-support queries, and each one
is implemented in all three query languages. Examples include:

\begin{itemize}[leftmargin=*]
    \item Top 10 products by total revenue.
    \item Monthly revenue trend.
    \item Top 10 customers by spend.
    \item Revenue by product category.
    \item Profit by region and state.
    \item Products with high sales but low profit.
    \item Products with low sales but high profit.
    \item Shipping mode with the highest profit ratio.
    \item Customers with the most orders.
    \item Months with the lowest profit.
\end{itemize}

The front-end exposes three views:

\begin{enumerate}[leftmargin=*]
    \item \textbf{Home page}: three cards for the three databases, plus a
          link to the side-by-side comparison view
          (Figure~\ref{fig:home}).
    \item \textbf{Explorer (\texttt{/explorer/<db>})}: lets the user pick a
          predefined query or write their own custom query (SQL, Cypher, or
          a MongoDB aggregation JSON). The result is shown as a table and,
          optionally, as a chart (Figure~\ref{fig:explorer}).
    \item \textbf{Comparison (\texttt{/compare})}: runs the same semantic
          query on all three databases and shows the three results
          side-by-side, with row counts and execution time for each.
\end{enumerate}

\begin{figure}[H]
\centering
% TODO: Screenshot of http://localhost:5000 landing page.
% Should show the three DB cards (PostgreSQL, Neo4j, MongoDB) and
% the "Side-by-Side Comparison" card.
\fbox{\parbox{0.9\textwidth}{\centering
\textit{Placeholder: screenshot of the umbrella home page at
\texttt{http://localhost:5000}, showing the three database cards
and the side-by-side comparison card.}
\vspace{5em}}}
\caption{Umbrella home page: one entry point to three backends.}
\label{fig:home}
\end{figure}

\begin{figure}[H]
\centering
% TODO: Screenshot of /explorer/postgres after running e.g.
% "Monthly revenue trend" with the chart toggled on.
\fbox{\parbox{0.9\textwidth}{\centering
\textit{Placeholder: screenshot of the PostgreSQL explorer running
``Monthly revenue trend'', with both the result table and the
chart visible.}
\vspace{6em}}}
\caption{Per-database explorer page with tabular and chart output.}
\label{fig:explorer}
\end{figure}

\subsection{User Scenarios}

\textbf{Scenario 1 -- Analyst looking at a trend.} Anna, a business analyst,
opens the Explorer for PostgreSQL, picks ``Monthly revenue trend'' and
toggles the chart view. She immediately sees how revenue has evolved over
the last two years.

\textbf{Scenario 2 -- Cross-database comparison.} Ben, a student, opens the
comparison view and runs ``Top 10 products by revenue''. He sees that all
three databases return the same ranking, but with different execution times,
and uses this to discuss the trade-offs of each model.

\textbf{Scenario 3 -- Custom query.} Clara, a power user, opens the Neo4j
explorer and writes her own Cypher query to find which shipping modes are
most common in California. She does not need to understand the underlying
graph model in detail because the front-end takes care of executing and
rendering the result.

\subsection{Dockerization}

Everything runs through \texttt{docker-compose}. The stack includes:

\begin{itemize}[leftmargin=*]
    \item \texttt{umbrella} -- the Flask web application.
    \item \texttt{postgres}, \texttt{mongo}, \texttt{neo4j} -- the three
          databases.
    \item \texttt{zookeeper}, \texttt{kafka}, \texttt{schema-registry},
          \texttt{ksqldb-server}, \texttt{ksqldb-cli} -- the streaming
          stack.
    \item \texttt{ksqldb-init} -- a one-shot container that applies the
          ksqlDB schema from \texttt{kafka/ksql/init.sql} on first run.
\end{itemize}

A single \texttt{docker compose up -d --build} starts the whole system.
This makes the project easy to demo and to hand over: no manual database
installs, no version mismatches.

\subsection{Data Streaming Functionality}

Every time the front-end runs a query (through \texttt{/api/run} or
\texttt{/api/compare}), the Flask app executes the query, returns the
result to the user as usual, \emph{and} publishes a small metadata event to
Kafka. The event contains:

\begin{lstlisting}[caption={Query execution event (JSON)}]
{
  "event_id": "uuid-...",
  "db_name": "postgres",
  "query_type": "predefined",
  "query_key": "top_products_by_revenue",
  "success": true,
  "row_count": 10,
  "elapsed_ms": 42.15,
  "executed_at": "2026-04-18T09:14:33Z"
}
\end{lstlisting}

On top of this stream, ksqlDB creates:

\begin{itemize}[leftmargin=*]
    \item \texttt{QUERY\_EXECUTIONS\_RAW} -- the raw event stream.
    \item \texttt{QUERY\_EXECUTION\_METRICS} -- a table of 5-minute windowed
          aggregates (execution count, average latency, latest timestamp)
          grouped by database name and query type.
    \item \texttt{FAILED\_QUERY\_EXECUTIONS} -- a filtered stream of failed
          executions, useful for monitoring.
\end{itemize}

The key design choice was to stream \emph{query execution metadata}, not
full query results. Metadata is small, uniform across all three backends,
and the most interesting for an operational dashboard (``how fast is the
warehouse answering us?''). It also keeps the primary user-facing path
unchanged: if Kafka is down, the user still gets the query result; only the
metric is lost.

% =====================================================
\section{Results}\label{sec:results}

\subsection{Working Front-End}

The umbrella web application is reachable at \url{http://localhost:5000}
after \texttt{docker compose up}. From the home page the user can:

\begin{itemize}[leftmargin=*]
    \item Open an explorer for any of the three databases.
    \item Run any of the ten predefined queries or write a custom one.
    \item Switch to the side-by-side comparison view.
    \item Toggle a chart for the returned result set.
\end{itemize}

\subsection{Same Question, Three Answers}

For every predefined query, running it through the comparison view returns
three result tables that agree on the business answer. For example, ``Top 10
products by revenue'' returns the same ranking on PostgreSQL, MongoDB, and
Neo4j, expressed through three very different query languages
(SQL, aggregation pipeline, Cypher).

This is the core deliverable of the project: the same semantic question is
answered by three structurally different systems, and the front-end makes it
easy to verify this.

\subsection{Streaming Metrics}

After running a handful of queries, the ksqlDB CLI shows live data:

\begin{itemize}[leftmargin=*]
    \item \texttt{SELECT * FROM QUERY\_EXECUTIONS\_RAW EMIT CHANGES;} prints
          one row per query executed.
    \item \texttt{SELECT * FROM QUERY\_EXECUTION\_METRICS EMIT CHANGES;}
          prints rolling 5-minute averages per database and query type.
    \item \texttt{SELECT * FROM FAILED\_QUERY\_EXECUTIONS EMIT CHANGES;}
          shows only failed executions (useful to catch syntax errors in
          custom queries).
\end{itemize}

This gives a lightweight operational view of the warehouse without any
extra monitoring tools.

\subsection{Reproducibility}

Because the whole stack is in Docker, the project can be reproduced from a
clean machine in one command. The \texttt{ksqldb-init} container applies the
ksqlDB schema automatically on first startup, so there is no manual setup
drift between runs.

% =====================================================
\section{Discussions}\label{sec:discussion}

Building the same warehouse three times gave us a better view of where
the boundaries really are between relational, document, and graph
systems. This section pulls together what we actually observed and
discusses why each model behaved the way it did, rather than simply
repeating what the textbooks say.

\subsection{Same Semantics, Three Different Worlds}

The central experiment of the project was to take ten decision-support
queries and express each of them in SQL, MongoDB aggregation pipelines,
and Cypher. Figure~\ref{fig:three-models} puts the three data
representations next to each other for a single sales fact.

\begin{figure}[H]
\centering
% TODO: Replace with a composite figure (3 panels).
% Panel (a) -- Star schema ERD (you can reuse your Supabase ERD screenshot).
% Panel (b) -- A single MongoDB sales document with embedded product,
%              customer, location, ship_mode and date sub-documents.
% Panel (c) -- Neo4j graph view: a (:Sale) node linked via labelled
%              relationships to (:Product), (:Customer), (:State),
%              (:ShipMode), (:Date) nodes.
\fbox{\parbox{0.95\textwidth}{\centering
\textit{Placeholder: three-panel figure showing the same sale as
(a) a row in \texttt{fact\_sales} plus dimension rows,
(b) a single embedded MongoDB document,
(c) a subgraph in Neo4j with labelled relationships.}
\vspace{6em}}}
\caption{The same business fact expressed in the relational, document and graph models.}
\label{fig:three-models}
\end{figure}

What became clear, once we looked at these side by side, is that the
three models are not interchangeable surface syntaxes. They encode
different assumptions about what is stable, what is joined, and what is
expensive. The star schema assumes that dimensions are shared and
normalised; the document model assumes that each fact carries its own
context; the graph model assumes that relationships between entities are
themselves the interesting data.

\subsection{Modelling Trade-offs}

Table~\ref{tab:compare} summarises the characteristics we observed while
implementing the warehouse three times. It is a refinement of the general
NoSQL trade-offs discussed in the
literature~\cite{cattell2011nosql,gessert2017nosql,grolinger2013cloud},
specialised to what a small data-warehousing project like ours actually
cares about.

\begin{table}[H]
\centering
\caption{Relational vs. NoSQL implementations for the same warehouse, as observed in this project.}
\label{tab:compare}
\begin{tabularx}{\textwidth}{l X X X}
\toprule
\textbf{Aspect} & \textbf{PostgreSQL (SQL)} & \textbf{MongoDB (Documents)} & \textbf{Neo4j (Graph)} \\
\midrule
Conceptual shape     & Fact + dimension tables                 & Fact document with embedded context & Labelled nodes and typed relationships \\
Modelling effort     & Moderate (normalise into star schema)   & Low (one document per sale)         & Moderate (decide nodes/relationships) \\
Query language       & Standard SQL                            & JSON aggregation pipeline           & Cypher (pattern matching)             \\
Joins                & Explicit via keys                       & Mostly avoided (embedded)           & Traversals via relationships          \\
Schema flexibility   & Fixed schema                            & Schema-on-read                      & Flexible, constraints optional        \\
Aggregation ergonomics & Excellent (\texttt{GROUP BY}, window fns) & Good but verbose pipelines       & Adequate but less natural than SQL    \\
Relationship queries & Possible but multi-join                 & Limited without \texttt{\$lookup}   & First-class (paths, traversals)       \\
Dimension updates    & One row in one table                    & Rewrite every embedded copy         & Update a single dimension node        \\
Strengths for DW     & BI-friendly, mature optimiser           & Fast single-document reads          & Natural for relationship-heavy analysis \\
Weaknesses for DW    & Many joins on wide stars                & Dimension drift across documents    & Aggregations less natural             \\
Typical learning curve & Low (if you know SQL)                 & Low-to-medium                       & Medium (new mental model)             \\
\bottomrule
\end{tabularx}
\end{table}

\subsection{Query Expressiveness and Readability}

For simple aggregations -- ``top 10 products by revenue'', ``monthly
revenue trend'' -- all three languages produce compact code of similar
length. PostgreSQL wins on readability for anyone who already knows SQL;
Cypher wins when the query is shaped like a path through the graph; and
MongoDB's pipeline is the most verbose but is also the closest to a step
by step data-flow description, which some team members preferred.

As soon as questions involve more than one level of aggregation or a
non-trivial relationship, the differences become more visible. For
example, a hypothetical query such as ``for each product, which customers
have bought it more than three times?'' would need a self-join in
PostgreSQL, a nested \texttt{\$group} stage in MongoDB, and would read
almost like a sentence in Cypher. None of the predefined queries in this
project go that far, but it is a good example of where the graph model
would start to pull ahead.

\subsection{A Note on Elapsed-Time Numbers}

The side-by-side comparison view (Figure~\ref{fig:compare-screenshot})
reports the elapsed time per backend for every query, and those numbers
are also the ones that flow through Kafka into the ksqlDB metrics. It
would be tempting to use them to rank the three databases, but we have
deliberately avoided doing so.

\begin{figure}[H]
\centering
% TODO: Screenshot of /compare after running e.g. "Top 10 products by revenue".
% Show the three result cards with row counts and elapsed_ms badges visible.
\fbox{\parbox{0.9\textwidth}{\centering
\textit{Placeholder: screenshot of the side-by-side comparison view
(\texttt{/compare}) after running ``Top 10 products by revenue'',
with three result cards (PostgreSQL, Neo4j, MongoDB) and their
row counts and \texttt{ms} timings visible.}
\vspace{6em}}}
\caption{The comparison view surfaces per-query latency as an operational signal, not as a benchmark.}
\label{fig:compare-screenshot}
\end{figure}

The reason is that all three databases in our setup are hosted as
managed cloud services (PostgreSQL via Supabase, Neo4j Aura, MongoDB
Atlas), and the Flask application talks to them over TLS from Docker.
What the elapsed-time numbers actually reflect is a mixture of three
effects that have little to do with the data model itself: whether the
Python client keeps a warm connection pool or reconnects per query, how
the managed service fronts requests (for example, Supabase routes
queries through a pgbouncer transaction-mode pooler on port
\texttt{6543}), and WAN round-trip latency between the client and each
cloud region.

For that reason we treat \texttt{elapsed\_ms} in this project as an
\emph{operational observability signal} -- useful for spotting slow
queries, trends, and failures through the ksqlDB metrics pipeline --
and not as a measurement of which engine is ``faster''. A fair
head-to-head comparison would require connection pooling on all three
clients, co-located databases, identical index coverage, and controlled
warm-up, none of which are in scope for this project.

\subsection{Value of the Streaming Layer}

Adding Kafka and ksqlDB was not strictly required by the assignment, but
in hindsight it changed the character of the project. Because every
query goes through an event (Figure~\ref{fig:ksql-metrics}), we get an
observability layer on top of a system that would otherwise be
black-box. It is striking how cheap this was to add: a single producer
call in the Flask app, a \texttt{ksqldb-init} container, and a few lines
of ksqlDB SQL.

\begin{figure}[H]
\centering
% TODO: Terminal screenshot from the ksqlDB CLI showing output of
% SELECT * FROM QUERY_EXECUTION_METRICS EMIT CHANGES LIMIT 5;
% after running a handful of queries across all three DBs.
\fbox{\parbox{0.9\textwidth}{\centering
\textit{Placeholder: ksqlDB CLI session showing a \texttt{SELECT * FROM
QUERY\_EXECUTION\_METRICS EMIT CHANGES} output, with per-database
rolling counts and average latencies updating in real time.}
\vspace{5em}}}
\caption{Real-time query-execution metrics derived by ksqlDB from the \texttt{dw.query.executions} topic.}
\label{fig:ksql-metrics}
\end{figure}

The design choice that paid off was streaming \emph{metadata}, not
results. Query results vary wildly in shape between backends (rows,
documents, records), so forcing them through a single stream would have
required a non-trivial translation layer. Metadata is uniform across all
three databases and is exactly what an operator of the warehouse wants
to see.

\subsection{When Would We Choose Which?}

Looking at the three implementations together, a practical rule of
thumb emerged for us:

\begin{itemize}[leftmargin=*]
    \item If the workload is classic BI on a well-understood domain,
          the \textbf{relational} warehouse is the default. The star
          schema is easy to reason about, SQL is familiar to analysts
          and BI tools, and the ecosystem around it is mature.
    \item If the data is naturally ``one event with a lot of context'',
          is evolving quickly, and is mostly read as-is, a
          \textbf{document} store like MongoDB removes a lot of
          modelling friction. Dimension churn is the main risk, since
          embedded fields drift out of sync when source data changes.
    \item If the interesting questions are \emph{about the
          relationships} -- recommendation, fraud, pathfinding, network
          analysis -- a \textbf{graph} database like Neo4j becomes the
          natural home. Cypher expresses these questions almost as
          sentences, which neither SQL nor aggregation pipelines can
          match.
\end{itemize}

Notice that none of these bullets are about raw speed. With this much
variation in hosting, client libraries, and workload mix, we felt that
picking a winner on latency alone would be misleading. In many real
organisations the answer also ends up being ``all of the above'', and an
umbrella front-end like the one we built is a reasonable way to hide
that from end users.

\subsection{Limitations and Threats to Validity}

We should be honest about what this project does and does not show.

\begin{itemize}[leftmargin=*]
    \item \textbf{Data volume.} Our dataset is small and fits comfortably
          in memory on every backend, so the project does not say
          anything useful about how the three models scale.
    \item \textbf{Not a benchmark.} The \texttt{elapsed\_ms} values shown
          in the UI and aggregated by ksqlDB are a user-facing
          operational signal, not a controlled measurement. They are
          dominated by client-side connection strategy and WAN latency
          to the managed cloud services (Supabase, Neo4j Aura, MongoDB
          Atlas) rather than by the engines themselves, and we
          explicitly do not use them to rank the databases.
    \item \textbf{Cross-backend equivalence.} We wrote the queries by
          hand in three languages. For the simple aggregations the
          equivalence is easy to argue for, but more complex queries
          would benefit from a shared intermediate representation to
          rule out accidental semantic drift.
    \item \textbf{No real ETL.} We treat the cleaned dataset as the
          input. A real deployment would need a nightly load job, a
          summary-table refresh job, and proper failure handling -- the
          project only describes these at a specification level.
    \item \textbf{Observational streaming only.} Kafka + ksqlDB watch
          queries but do not drive decisions. A natural next step would
          be to close the loop, e.g. by using rolling latency to auto-
          suggest which backend to run a given query on.
    \item \textbf{No security story.} Authentication, authorisation and
          audit logging are out of scope for a teaching prototype and
          would need to be added before anything like this could face
          real users.
\end{itemize}

% =====================================================
\section{Conclusions}\label{sec:conclusion}

We built a data warehouse for a simplified e-commerce business and
implemented it three times in three very different database technologies: a
classic relational star schema, a document model, and a graph model. A
single Flask-based umbrella front-end lets users run decision-support
queries on any of the backends and compare the results side-by-side. A
Kafka + ksqlDB pipeline adds a lightweight layer of observability on top of
the whole system, without changing the primary user experience.

The main takeaway is that data-warehousing concepts -- facts, dimensions,
aggregations -- are not tied to any particular database technology. They can
be expressed naturally in relational, document, and graph systems, though
each representation has trade-offs. PostgreSQL remains the easiest target
for classic BI-style analytics; MongoDB is attractive when documents map
naturally to ``one event with context''; and Neo4j is most interesting when
the questions are really about relationships.

Future work could include a proper ETL pipeline into all three stores, more
advanced ksqlDB derivations (such as anomaly detection on query latency),
and richer graph-shaped questions in Neo4j that go beyond what the
relational star schema can express comfortably.

% =====================================================
\section{Use of AI}\label{sec:ai}

We used AI assistants (large language model based) during the project as a
productivity tool. Concretely, AI was used for:

\begin{itemize}[leftmargin=*]
    \item Generating initial drafts of boilerplate (Flask routes, Docker
          Compose skeleton) that we then reviewed and adapted.
    \item Translating our PostgreSQL queries into equivalent MongoDB
          aggregation pipelines and Cypher, which we then tested against
          the actual data.
    \item Explaining unfamiliar error messages, especially around ksqlDB
          startup and Kafka listeners.
    \item Improving the wording and structure of this report.
\end{itemize}

All code was read, tested, and modified by us. All decisions about the data
model, the query set, and the system architecture were made by the team.
The AI was used as a helper, not as an author.

% =====================================================
\newpage
\printbibliography[title={References}]
\addcontentsline{toc}{section}{References}

% =====================================================
\newpage
\appendix
\section{Datasheet}\label{app:datasheet}

\subsection*{A.1 Source data}

\begin{itemize}[leftmargin=*]
    \item Domain: e-commerce (orders, products, customers, locations,
          shipping modes).
    \item Origin: synthesised / simulated dataset inspired by public
          retail datasets.
    \item Size: on the order of $10^4$--$10^5$ sales records (sufficient
          for all demo queries to run in well under a second).
\end{itemize}

\subsection*{A.2 Fields used in the warehouse}

\begin{itemize}[leftmargin=*]
    \item Sale measures: \texttt{sales\_amount}, \texttt{quantity},
          \texttt{profit\_amount}.
    \item Product: \texttt{product\_name}, \texttt{category\_name}.
    \item Customer: \texttt{customer\_name}.
    \item Location: \texttt{region}, \texttt{state}.
    \item Shipping: \texttt{ship\_mode\_name}.
    \item Date: \texttt{year}, \texttt{month}, \texttt{month\_name}.
\end{itemize}

\subsection*{A.3 Environment}

\begin{itemize}[leftmargin=*]
    \item PostgreSQL 16, MongoDB 7, Neo4j 5.22.
    \item Confluent Platform 7.6.1 (Kafka, Schema Registry, ksqlDB).
    \item Python 3.11, Flask 3.0, \texttt{psycopg2-binary}, \texttt{pymongo},
          official Neo4j Python driver, \texttt{kafka-python}.
    \item Docker and \texttt{docker-compose}.
\end{itemize}

\subsection*{A.4 Key endpoints}

\begin{itemize}[leftmargin=*]
    \item Web UI: \url{http://localhost:5000}
    \item ksqlDB REST API: \url{http://localhost:8088}
    \item Schema Registry: \url{http://localhost:8081}
\end{itemize}

\end{document}
